\documentclass[a4paper,12pt]{article}
\usepackage[T2A]{fontenc}
\usepackage[utf8]{inputenc}
\usepackage[russian]{babel}
\usepackage{amsmath,amssymb}
\usepackage{geometry,booktabs,tabularx,array,enumitem,fancyhdr,setspace,xcolor,tikz,microtype,titlesec}
\usetikzlibrary{arrows.meta,positioning,shapes.geometric,calc}
\geometry{left=20mm,right=20mm,top=20mm,bottom=22mm,headheight=25pt,headsep=5mm}
\definecolor{accent}{HTML}{008080}
\definecolor{darkaccent}{HTML}{004D4D}
\definecolor{textgray}{HTML}{333333}
\color{textgray}
\setlength{\parindent}{0pt}
\setlength{\parskip}{5pt}
\setstretch{1.22}
\pagestyle{fancy}
\fancyhf{}
\fancyhead[L]{\small Движение и производительность}
\fancyhead[R]{\small Ярослав Гаврилов}
\fancyfoot[C]{\small\thepage}
\renewcommand{\headrulewidth}{0.4pt}
\renewcommand{\headrule}{\hbox to\headwidth{\color{accent}\leaders\hrule height\headrulewidth\hfill}}
\titleformat{\section}{\normalfont\Large\color{darkaccent}}{\thesection.}{0.5em}{}[\vspace{2mm}{\color{accent}\titlerule[0.4pt]}]
\titlespacing*{\section}{0pt}{8pt}{9pt}
\newcommand{\example}[1]{\par\medskip{\color{darkaccent}\textbf{#1}}\par\smallskip}
\newcommand{\unit}[1]{\,\text{#1}}
\newcommand{\checkitem}{\item[\(\square\)]}
\setlist[itemize]{leftmargin=6mm,itemsep=2pt,topsep=3pt}
\setlist[enumerate]{leftmargin=8mm,itemsep=8pt,topsep=4pt}
\renewcommand{\arraystretch}{1.32}
\begin{document}
\begin{titlepage}
\centering
\vspace*{28mm}
{\Large\color{darkaccent} МАТЕМАТИКА\par}
\vspace{17mm}
{\Huge Чек-лист\par}
\vspace{11mm}
{\LARGE\color{accent} Текстовые задачи на движение\\[2mm] и производительность\par}
\vspace{18mm}
{\large Подготовка к экзамену по математике\par}
\vfill
{\large Ярослав Гаврилов\par}
\vspace{10mm}
{\large 26 ноября 2026 г.\par}
\vspace{17mm}
{\small Лёвин Артём Александрович\\эксклюзивно для Ярослава Гаврилова\par}
\vspace{17mm}
\begin{tikzpicture}
\draw[accent,line width=.65pt] (0,0) .. controls (2,0.8) and (4,-0.8) .. (6,0) .. controls (8,0.8) and (10,-0.8) .. (12,0);
\end{tikzpicture}
\vspace{4mm}
\end{titlepage}
\setcounter{page}{1}
\section{Единицы измерения и движение поездов}
Главная идея: сначала выберите согласованные единицы измерения, затем составляйте уравнение. Скорость определяет, в каких единицах удобно выражать расстояние и время.
\[
 S=vt,\qquad v=\frac{S}{t},\qquad t=\frac{S}{v}.
\]
\begin{center}
\begin{tabularx}{\linewidth}{@{}lXX@{}}
\toprule
Скорость & Расстояние & Время\\\midrule
км/ч & километры & часы\\
м/с & метры & секунды\\
м/мин & метры & минуты\\\bottomrule
\end{tabularx}
\end{center}
\[
1\unit{км}=1000\unit{м},\qquad 1\unit{ч}=60\unit{мин}=3600\unit{с}.
\]
Из метров в километры делите на 1000, из минут в часы --- на 60. Из часов в секунды умножайте на 3600. При переводе скорости из км/ч в м/с делите на $3{,}6$.

\example{Полное расхождение поездов}
При движении по параллельным путям от встречи локомотивов до расхождения последних вагонов относительное расстояние равно сумме длин поездов.
\begin{center}
\begin{tikzpicture}[x=1cm,y=1cm,>=Stealth]
\draw[accent,thick,rounded corners=1pt] (-4,0) rectangle (0,.48);
\draw[darkaccent,thick,rounded corners=1pt] (0,-.8) rectangle (5,-.32);
\draw[densely dashed,gray] (0,.7)--(0,-1.02);
\draw[->,accent,thick] (-2.8,.65)--(-1,.65) node[midway,above] {$v_1$};
\draw[->,darkaccent,thick] (3,-.13)--(1.2,-.13) node[midway,above] {$v_2$};
\node at (-2,.24) {$L_1$}; \node at (2.5,-.56) {$L_2$};
\node[below,font=\small] at (0,-1.05) {встреча локомотивов};
\end{tikzpicture}
\end{center}
\[
 t=\frac{L_1+L_2}{v_1+v_2}\quad\text{(навстречу)},\qquad
 t=\frac{L_1+L_2}{v_{\text{быстр}}-v_{\text{медл}}}\quad\text{(обгон)}.
\]
При обгоне время отсчитывают от момента, когда передняя часть быстрого поезда поравнялась с задней частью медленного, до полного опережения.

\example{Пример 1. Поезда длиной 240 и 360 м, скорости 54 и 90 км/ч}
Поезда движутся навстречу. Требуется найти время в секундах.
\[
 t=\frac{0{,}24+0{,}36}{54+90}=\frac{0{,}6}{144}=\frac{1}{240}\unit{ч},
 \qquad \frac{3600}{240}=15\unit{с}.
\]
\textbf{Ответ:} 15 секунд. Не подставляйте метры в формулу со скоростью в км/ч.
\clearpage
\section{Производительность, время и результат}
Производительность --- это объём работы за \textbf{одну} единицу времени. Обозначим её $p$, время --- $t$, результат (объём выполненной работы) --- $A$.
\begin{center}
\begin{tikzpicture}[scale=1]
\draw[accent,thick] (-1.9,0)--(0,2.1)--(1.9,0)--cycle;
\draw[accent] (-1.05,.94)--(1.05,.94) (0,0)--(0,.94);
\node at (0,1.3) {$A$}; \node at (-.58,.36) {$p$}; \node at (.58,.36) {$t$};
\end{tikzpicture}
\hspace{12mm}
\begin{minipage}{.49\linewidth}
\[ A=pt,\qquad p=\frac{A}{t},\qquad t=\frac{A}{p}. \]
В рассматриваемых задачах $p>0$, $t>0$.
\end{minipage}
\end{center}
Если мастер делает заказ за 3 часа, его производительность равна $1/3$ заказа в час. Если труба пропускает 7 л за 4 минуты, то $p=7/4$ л/мин, а не 7 и не 28.

\example{Как заполнять таблицу}
Строки обозначают участников или режимы работы; столбцы --- производительность, время и результат. Впишите известные величины и неизвестную $x$, затем заполните оставшиеся клетки по формулам. Единицы каждого столбца должны совпадать.
\section{Совместная работа двух исполнителей}
При постоянных производительностях и одновременной независимой работе над общим результатом складываются производительности, а не времена:
\[p_{\text{общ}}=p_1+p_2,\qquad t_{\text{общ}}=\frac{A}{p_1+p_2}.\]
\example{Пример 2. Один заказ: первый мастер выполняет его за 3 ч, второй --- за 6 ч}
\begin{center}
\begin{tabularx}{\linewidth}{@{}Xccc@{}}
\toprule
Исполнитель & $p$, заказ/ч & $t$, ч & $A$, заказ\\\midrule
Первый & $1/3$ & 3 & 1\\
Второй & $1/6$ & 6 & 1\\
Вместе & $1/x$ & $x$ & 1\\\bottomrule
\end{tabularx}
\end{center}
\[\frac1x=\frac13+\frac16=\frac12\quad\Longrightarrow\quad x=2\unit{ч}.\]
Если объём неизвестен и спрашивают только время выполнения одной и той же работы, удобно принять весь заказ за единицу. Если объём известен и важен для условия (литры, детали), используйте его реальное значение.
\clearpage
\example{Пример 3. Две трубы наполняют сосуд}
Первая труба подаёт 7 л за 4 минуты, вторая --- 7 л за 5 минут. Требуется наполнить сосуд объёмом 63 л при одновременной работе труб.
\begin{center}
\begin{tabularx}{\linewidth}{@{}Xccc@{}}
\toprule
Труба & $p$, л/мин & $t$, мин & $A$, л\\\midrule
Первая & $7/4$ & 4 & 7\\
Вторая & $7/5$ & 5 & 7\\
Вместе & $63/x$ & $x$ & 63\\\bottomrule
\end{tabularx}
\end{center}
\[
\frac{63}{x}=\frac74+\frac75
\quad\Longrightarrow\quad \frac9x=\frac14+\frac15=\frac9{20}
\quad\Longrightarrow\quad x=20\unit{мин}.
\]
Сокращайте общий числовой множитель до приведения дробей к общему знаменателю. Проверка: за 20 минут трубы подадут $35+28=63$ л.
\section{Разность времени: одинаковый объём работы}
Если один исполнитель справляется быстрее другого, то из большего времени вычитайте меньшее. При одинаковом объёме работы большая производительность соответствует меньшему времени.
\example{Пример 4. Петя и Ваня решают один и тот же тест}
Петя отвечает на 6 вопросов в час, Ваня --- на 7. Петя заканчивает на 20 минут позже. Найдём число вопросов $x$.
\begin{center}
\begin{tabularx}{\linewidth}{@{}Xccc@{}}
\toprule
Исполнитель & $p$, вопр./ч & $t$, ч & $A$, вопр.\\\midrule
Петя & 6 & $x/6$ & $x$\\
Ваня & 7 & $x/7$ & $x$\\\bottomrule
\end{tabularx}
\end{center}
Время в часах, поэтому $20$ минут $=1/3$ часа. При $x>0$:
\[
 \frac{x}{6}-\frac{x}{7}=\frac13
 \quad\Longrightarrow\quad \frac{x}{42}=\frac13
 \quad\Longrightarrow\quad x=14.
\]
Проверка: Петя затратит $7/3$ ч, Ваня --- 2 ч. Разница равна 20 минутам. Ответ должен быть целым положительным числом.
\clearpage
\section{План и факт. Разность производительностей}
\example{Пример 5. Выполнение плана досрочно}
Нужно прочитать 300 страниц. Если читать на 5 страниц в день больше, чем планировалось, работа закончится на 3 дня раньше. Пусть $x$ --- плановое число страниц в день, $x>0$.
\begin{center}
\begin{tabularx}{\linewidth}{@{}Xccc@{}}
\toprule
Режим & $p$, стр./день & $t$, дни & $A$, стр.\\\midrule
По плану & $x$ & $300/x$ & 300\\
Фактически & $x+5$ & $300/(x+5)$ & 300\\\bottomrule
\end{tabularx}
\end{center}
\[
 \frac{300}{x}-\frac{300}{x+5}=3
 \quad\Longrightarrow\quad x(x+5)=500
 \quad\Longrightarrow\quad x^2+5x-500=0.
\]
Корни $20$ и $-25$. Отрицательный корень не подходит. Проверка: $300/20=15$ дней, $300/25=12$ дней. Ответ: 20 страниц в день по плану.
\example{Пример 6. Одинаковое время, разные объёмы}
Первая труба наполняет бак объёмом 770 л, вторая --- бак объёмом 830 л за такое же время. Одна труба подаёт на 6 л/мин больше другой. Какая из них производительнее?

При одинаковом положительном времени большему объёму соответствует большая производительность. Значит, у первой трубы она равна $x$, у второй --- $x+6$.
\begin{center}
\begin{tabularx}{\linewidth}{@{}Xccc@{}}
\toprule
Труба & $p$, л/мин & $t$, мин & $A$, л\\\midrule
Первая & $x$ & $770/x$ & 770\\
Вторая & $x+6$ & $830/(x+6)$ & 830\\\bottomrule
\end{tabularx}
\end{center}
\[
 \frac{770}{x}=\frac{830}{x+6}\quad(x>0),\qquad
 77(x+6)=83x\quad\Longrightarrow\quad x=77.
\]
Ответ: 77 и 83 л/мин соответственно. Обе трубы работают по 10 минут. Слово «одна» не указывает, какая именно труба работает быстрее: это нужно вывести из остальных данных.
\clearpage
\section{Последовательная работа: откачивание и заполнение}
Если процессы идут один за другим, складывайте их времена. Не путайте суммарное время с временем отдельного процесса. При одинаковом объёме меньшая производительность означает большее время.
\example{Пример 7. Бассейн объёмом 30 м$^3$}
Сначала из полного бассейна откачали всю воду, затем снова заполнили его. На обе операции ушло 8 часов. При заполнении насос перекачивает на 4 м$^3$/ч меньше, чем при откачивании. Найдём время заполнения.
\begin{center}
\begin{tabularx}{\linewidth}{@{}Xccc@{}}
\toprule
Операция & $p$, м$^3$/ч & $t$, ч & $A$, м$^3$\\\midrule
Откачивание & $30/(8-x)$ & $8-x$ & 30\\
Заполнение & $30/x$ & $x$ & 30\\\bottomrule
\end{tabularx}
\end{center}
Пусть $x$ --- время заполнения. Тогда $0<x<8$, а откачивание длится $8-x$ часов. Так как заполнение медленнее, $x>8-x$, то есть $x>4$.
\[
 \frac{30}{8-x}-\frac{30}{x}=4.
\]
Умножаем обе части на $x(8-x)$, ненулевое на допустимом интервале:
\begin{align*}
 30x-30(8-x)&=4x(8-x),\\
 60x-240&=32x-4x^2,\\
 x^2+7x-60&=0,\\
 (x-5)(x+12)&=0.
\end{align*}
Подходит только $x=5$. Проверка: на откачивание ушло $8-5=3$ часа, производительности равны $30/3=10$ и $30/5=6$ м$^3$/ч; разность равна 4.

\textbf{Ответ:} 5 часов.
\example{Что проверять в условии}
\begin{itemize}
\checkitem «На всё ушло 8 часов» означает сумму времён, а не время откачивания.
\checkitem «При заполнении на 4 меньше» относится к производительности, а не ко времени.
\checkitem В разности первой стоит большая величина: здесь производительность откачивания.
\checkitem При дробно-рациональном уравнении исключите нулевые знаменатели и отрицательные времена.
\end{itemize}
\clearpage
\section{Совместная работа трёх исполнителей}
Если известны времена выполнения одной работы тремя парами, не обязательно находить производительность каждого участника отдельно. Достаточно сразу найти сумму всех трёх производительностей.
\example{Пример 8. Игорь, Паша и Володя}
Игорь с Пашей выполняют работу за 3 ч, Паша с Володей --- за 6 ч, Володя с Игорем --- за 4 ч. За сколько минут они выполнят эту работу втроём?
\begin{center}
\begin{tabularx}{\linewidth}{@{}Xccc@{}}
\toprule
Участники & $p$, работа/ч & $t$, ч & $A$, работа\\\midrule
Игорь и Паша & $1/3$ & 3 & 1\\
Паша и Володя & $1/6$ & 6 & 1\\
Володя и Игорь & $1/4$ & 4 & 1\\
Все трое & $1/x$ & $x$ & 1\\\bottomrule
\end{tabularx}
\end{center}
Обозначим индивидуальные производительности $p_1$, $p_2$, $p_3$. Тогда
\[
\begin{cases}
 p_1+p_2=1/3,\\
 p_2+p_3=1/6,\\
 p_3+p_1=1/4.
\end{cases}
\]
Сложим левые и правые части трёх равенств. Каждая производительность встретится ровно два раза:
\[
 2(p_1+p_2+p_3)=\frac13+\frac16+\frac14=\frac34.
\]
\[
 p_1+p_2+p_3=\frac38,\qquad x=\frac{1}{3/8}=\frac83\unit{ч}.
\]
Переводим в минуты: $\frac83\cdot60=160$. Ответ: 160 минут.

\example{Общий приём для трёх пар}
При временах $t_{12},t_{23},t_{31}$, выраженных в одинаковых единицах, и одной и той же работе:
\[
 t_{\text{втроём}}=\frac{2}{\dfrac{1}{t_{12}}+\dfrac{1}{t_{23}}+\dfrac{1}{t_{31}}}.
\]
Время втроём должно быть меньше времени любой пары.\par
\textbf{Самопроверка:} если три пары моют одно окно за 12, 20 и 15 минут, то $1/12+1/20+1/15=1/5$, общая производительность равна $1/10$ окна в минуту. Ответ: 10 минут.
\clearpage
\section*{Алгоритм: задачи на производительность}
\begin{center}
\begin{tikzpicture}[
 node distance=8mm and 7mm,
 action/.style={draw=accent,rounded corners=2pt,align=center,text width=10.6cm,minimum height=11mm,inner sep=5pt,font=\small},
 branch/.style={action,text width=4.8cm,minimum height=25mm},
 term/.style={draw=darkaccent,ellipse,align=center,inner sep=6pt,font=\small},
 decision/.style={diamond,draw=darkaccent,aspect=3,align=center,inner sep=2pt,font=\small},
 arrow/.style={-{Stealth[length=2.2mm]},draw=darkaccent,line width=.65pt}]
\node[term] (start) {Прочитайте условие};
\node[action,below=of start] (units) {Согласуйте единицы измерения};
\node[action,below=of units] (table) {Введите $x$ и ограничения. Заполните\\ таблицу $p$, $t$, $A$ по формулам};
\node[decision,below=of table] (joint) {Работают одновременно?};
\node[branch,below left=12mm and -12mm of joint] (yes) {Сложите производительности. Для трёх пар сложите три равенства и разделите на 2};
\node[branch,below right=12mm and -12mm of joint] (no) {Используйте связь из текста: равные времена, разность величин или сумму времён};
\node[action,below=43mm of joint] (equation) {Составьте уравнение и решите его с учётом ограничений};
\node[decision,below=of equation] (verify) {Есть допустимый корень?};
\node[action,below=of verify] (answer) {Переведите результат в требуемые единицы};
\node[term,below=of answer] (end) {Запишите ответ};
\draw[arrow] (start)--(units); \draw[arrow] (units)--(table); \draw[arrow] (table)--(joint);
\draw[arrow] (joint.west)-| node[pos=.25,above,font=\small]{да} (yes.north);
\draw[arrow] (joint.east)-| node[pos=.25,above,font=\small]{нет} (no.north);
\coordinate (merge) at ($(equation.north)+(0,5mm)$);
\draw[draw=darkaccent,line width=.65pt] (yes.south)|-(merge);
\draw[draw=darkaccent,line width=.65pt] (no.south)|-(merge);
\draw[arrow] (merge)--(equation.north);
\draw[arrow] (equation)--(verify);
\draw[arrow] (verify)--node[right,font=\small]{да}(answer);
\draw[arrow] (answer)--(end);
\draw[arrow] (verify.east)--++(3.35,0) node[midway,above,font=\small]{нет} |- (table.east);
\end{tikzpicture}
\end{center}
\clearpage
\section*{Тренировочный блок}
В каждой задаче на производительность составьте таблицу, укажите ограничения, запишите уравнение и решите его. В задачах на движение сначала согласуйте единицы измерения. Все скорости и производительности постоянны в пределах описанного режима работы.
\example{Базовый и средний уровень}
\begin{enumerate}
\item По двум параллельным железнодорожным путям навстречу друг другу движутся два поезда. Длина первого поезда составляет 180 м, а второго --- 270 м. Скорость первого поезда равна 72 км/ч, второго --- 36 км/ч. Сколько секунд пройдёт с момента встречи их локомотивов до момента, когда последние вагоны полностью разойдутся?
\item Два поезда движутся по параллельным путям в одном направлении. Быстрый поезд длиной 150 м движется со скоростью 90 км/ч, а медленный поезд длиной 250 м --- со скоростью 54 км/ч. Отсчёт времени начинается, когда локомотив быстрого поезда поравнялся с последним вагоном медленного. Через сколько секунд быстрый поезд полностью опередит медленный?
\item Двум мастерам поручено изготовить один и тот же комплект деталей. Первый мастер способен выполнить весь заказ самостоятельно за 4 часа, а второй --- за 12 часов. Мастера приступают к работе одновременно и могут работать независимо друг от друга. За сколько часов они выполнят один комплект деталей, если будут работать вместе?
\item В пустую ёмкость объёмом 75 л воду подают через две трубы. Через первую трубу за каждые 4 минуты поступает 8 л воды, через вторую за каждые 5 минут --- 15 л. Обе трубы открывают одновременно и не закрывают до полного заполнения ёмкости. За сколько минут ёмкость наполнится?
\item Два школьника независимо друг от друга выполняют один и тот же тест, содержащий одинаковое для обоих число вопросов. Первый школьник отвечает на 8 вопросов в час, второй --- на 12 вопросов в час. Они начали одновременно, но первый завершил тест на 30 минут позже второго. Сколько вопросов содержалось в тесте?
\end{enumerate}
\clearpage
\section*{Тренировочный блок (продолжение)}
\example{Уровень выше среднего}
\begin{enumerate}[start=6]
\item Ученик планировал прочитать книгу объёмом 240 страниц, ежедневно читая одно и то же количество страниц. Он решил читать каждый день на 4 страницы больше, чем первоначально запланировал. Благодаря этому он закончил чтение на 2 дня раньше намеченного срока. Сколько страниц в день ученик собирался читать по первоначальному плану?
\item Первая труба наполняет первый пустой бак, а вторая --- второй пустой бак. Объём первого бака составляет 420 л, а второго --- 540 л. Одна из труб пропускает за минуту на 4 л воды больше другой. Если обе трубы открыли одновременно, то оба бака оказались полностью заполненными в один и тот же момент. Найдите производительность каждой трубы в литрах в минуту.
\item Рабочие обслуживают бассейн объёмом 24 м$^3$. Сначала насос полностью откачивает воду из заполненного бассейна, после чего тот же насос вновь наполняет его до исходного уровня. На две операции вместе требуется 10 часов, без учёта времени переключения насоса. При заполнении насос перекачивает на 2 м$^3$ воды в час меньше, чем при откачивании. Сколько часов занимает заполнение бассейна?
\end{enumerate}
\example{Сложный уровень}
\begin{enumerate}[start=9]
\item Три сотрудника могут выполнить один и тот же заказ, работая попарно или втроём. Первый и второй сотрудники справляются с заказом за 18 минут, второй и третий --- за 24 минуты, а первый и третий --- за 36 минут. Каждый сотрудник сохраняет свою производительность независимо от состава группы. За сколько минут они выполнят один заказ, если приступят к работе одновременно втроём?
\item Три бригады могут подготовить один и тот же участок к ремонту. Первая и вторая бригады при совместной работе справляются за 1 час, вторая и третья --- за 80 минут, а третья и первая --- за 2 часа. Производительность каждой бригады постоянна, объём работы во всех трёх случаях одинаков. За какое время участок будет подготовлен, если все три бригады начнут работать одновременно? Выразите ответ в минутах и секундах.
\end{enumerate}
\clearpage
\section*{Ответы}
\begin{center}
\begin{tabularx}{.88\linewidth}{@{}rX@{}}
\toprule
№ & Ответ\\\midrule
1 & 15 с\\
2 & 40 с\\
3 & 3 ч\\
4 & 15 мин\\
5 & 12 вопросов\\
6 & 20 страниц в день\\
7 & первая труба --- 14 л/мин; вторая --- 18 л/мин\\
8 & 6 ч\\
9 & 16 мин\\
10 & 53 мин 20 с\\\bottomrule
\end{tabularx}
\end{center}
\vspace{6mm}
\begin{center}{\large\color{darkaccent} Итоговая памятка}\end{center}
\begin{tabularx}{\linewidth}{@{}p{.37\linewidth}X@{}}
\toprule
Сигнал в условии & Что сделать\\\midrule
Км/ч и метры & Перевести расстояние в км или скорость в м/с.\\
Результат за несколько минут & Разделить результат на время, чтобы найти производительность.\\
Работают одновременно & Сложить производительности.\\
На несколько дней раньше & Вычесть меньшее время из большего.\\
За одинаковое время & Приравнять выражения для времени.\\
На всё ушло заданное время & Сложить времена последовательных операций.\\
Три пары, найти время втроём & Сложить три равенства для парных производительностей и разделить на 2.\\\bottomrule
\end{tabularx}
\medskip
Перед записью ответа проверьте единицы, допустимость корня и соответствие результата тому, что спрашивается в условии.
\end{document}
